% RA-L draft (6 pages max incl. references). Build: tectonic main.tex
\documentclass[letterpaper, journal, twoside]{IEEEtran}

\usepackage{amsmath,amssymb}
\usepackage{graphicx}
\usepackage{booktabs}
\usepackage{xcolor}
\usepackage[hidelinks]{hyperref}
\usepackage{cite}

\newcommand{\todo}[1]{\textcolor{red}{[TODO: #1]}}
\newcommand{\pizerofive}{$\pi_{0.5}$}

\title{Can a VLM Grade the Robot?\\ False Positives, Not Accuracy, Decide Whether Self-Improvement Works}

\author{\todo{Authors}%
\thanks{\todo{Affiliation, Georgia Institute of Technology.}}}

\markboth{IEEE Robotics and Automation Letters. Preprint Version. Month, Year}%
{Authors: Can a VLM Grade the Robot?}

\begin{document}
\maketitle

\begin{abstract}
In order for robot policies to improve from their own experience, something must decide which attempts succeeded.
In simulation the environment provides this signal, but in the real world it does not, and recent work proposes to
use vision-language models (VLMs) as graders in its place. It remains an open question whether such graders are good
enough, and what ``good enough'' means for the policy that is trained on their decisions. We study this question with
the \pizerofive{} policy in LIBERO, where the simulator's success check provides ground truth for both grading and
evaluation. We measure how often current VLM graders pass failed attempts, compare them with a small probe trained on
in-domain labels, and fine-tune the policy on datasets with controlled grader errors. Our experiments show that
off-the-shelf graders pass a large fraction of failures, nearly all of them near-misses, and that these false positives,
rather than rejected successes, are what erase the benefit of self-improvement. A real grader's effect on the policy
matches that of random errors at the same rates, which implies that the quantity to measure and control is the
grader's precision on the policy's own rollouts, rather than its accuracy. \todo{Final numbers with two seeds; strict-threshold result.}
\end{abstract}

\begin{IEEEkeywords}
\todo{Learning from experience, reinforcement learning, vision-language models, reward modeling.}
\end{IEEEkeywords}

\section{Introduction}

\IEEEPARstart{R}{obot} policies trained on demonstrations are limited by the demonstrations they receive.
Vision-language-action (VLA) models such as \pizerofive{}~\cite{pi05} can perform a wide range of manipulation tasks,
but no demonstration set covers every task, object, and scene that a robot will encounter, and collecting more
demonstrations is expensive. A natural way past this limit is to let the robot learn from its own attempts: the robot
tries a task, each attempt is labelled as a success or a failure, and the policy is fine-tuned on its successes, either
by filtered imitation or by reinforcement learning~\cite{sefm,pistar06,simplevlarl}. This kind of self-improvement
requires no new demonstrations, but it does require a judge. In simulation the judge is free, because the simulator can
check whether the goal holds. On a real robot there is no such check, and the self-improvement loop can only be as good
as whatever labels the attempts.

A person who watches a robot can usually tell whether it put the bowl on the plate, and it is reasonable to ask whether
a VLM can do the same. Early work used image-text similarity as a reward~\cite{vlmrm,roboclip} or asked a VLM
directly whether a task succeeded~\cite{successvqa}, and more recent robot reward models are trained on robot video,
including failures, to predict progress or success~\cite{gvl,roboreward,robometer}. These graders are typically
evaluated by their accuracy on labelled episodes. However, a self-improvement loop does not consume accuracy: it consumes
the set of episodes that the grader passes, and the two kinds of grader error enter that set differently. A false
positive, a failure that the grader passes, places a failed trajectory in the training set as if it were a
demonstration of success. A false negative, a success that the grader rejects, only removes a good trajectory. We might
therefore hypothesize that the two errors have very different costs, and that a grader's usefulness for
self-improvement is determined primarily by its false positives. Prior work points the same way: resampling against an
imperfect verifier is capped by its false positives~\cite{flaws}, injected false-positive rewards hurt reinforcement
learning more than missing ones~\cite{huang2025darkside}, and robot reward models that over-score failures train poor
policies~\cite{roboref,tian2026bad}. However, no work has isolated the two errors in the data used to fine-tune a robot
policy on its own rollouts. This setting is sharpened by a property of robot failures: many of them are near-misses,
such as a drawer opened but not the right one, that look almost exactly like successes, and current graders tend to
pass them~\cite{failbench,tian2026bad}.

In this paper, we study the effect of grader errors on self-improvement directly. We use \pizerofive{} in
LIBERO~\cite{libero}, where the simulator's success check provides ground truth, so that we can both measure the
errors of a grader and control them. We first measure how current VLM graders perform on the policy's own rollouts and
characterize the failures they pass. We then ask how much in-domain data is needed to build a better grader, by
comparing the graders with a linear probe trained on a small number of labelled rollouts. Finally, we fine-tune the
policy on datasets constructed from the same pool of rollouts with controlled rates of false positives and false
negatives, on the selection of a real VLM grader, and on random errors matched to that grader's rates. This last
comparison separates two explanations for a grader's effect on the policy: whether it depends only on how many errors
the grader makes, or also on which episodes it gets wrong.

Our central finding is that false positives, not false negatives, determine whether self-improvement succeeds.
Fine-tuning on true successes improves the policy, adding failures to the training set erases this improvement, and
removing a large fraction of the successes costs little. A real VLM grader's effect on the policy is matched by random
errors at its rates, so its usefulness can be predicted from quantities that are measurable on a small labelled set.
\todo{One sentence on the strict-threshold result and the share of failures in the kept set.} Together, these results
suggest that a grader for self-improvement should be selected and thresholded by its precision on the policy's own
rollouts, rather than by its accuracy.

This finding is not only a limitation. A policy that improves itself from its own experience, judged by a model rather
than a person, is one of the most appealing goals in robot learning, because it would let a robot keep getting better
without anyone collecting new demonstrations. Our results identify what such a loop needs. Because rejected successes
cost little, a grader does not need to recognize every success: it can be strict, provided that what it passes is
clean. In the second part of the paper, we therefore ask how to obtain a grader with high precision cheaply, and test
three ways of doing so: adapting the grader to the domain, training a small probe on a few hundred labelled rollouts,
and ranking episodes within each task rather than thresholding scores across tasks. \todo{Result: which fixes reach
low $\rho$ and recover Clean's gain.} Finally, the share of failures in the kept set falls as the policy's success rate
rises, so a loop that begins above a threshold makes its own grading easier with each round. We discuss what this
implies for recursive self-improvement.

\section{Related Work}

\textbf{VLMs as rewards and success detectors.} A number of prior works have used pre-trained vision-language models to
judge robot behaviour, either as a dense reward from image-text similarity~\cite{vlmrm,roboclip}, as a success detector
queried in natural language~\cite{successvqa}, or as an estimator of task progress from video~\cite{gvl}. More recent
robot reward models are trained on large collections of robot video that include failed and partially successful
attempts~\cite{roboreward,robometer}. These works, and recent audits of grader reliability~\cite{failbench,tian2026bad},
evaluate graders by their agreement with human or oracle labels. The closest work to ours relates reward-model errors to
the policies they train: RoboRef finds that offline accuracy is a poor predictor of robot RL success and that
over-scored failures are the main cause~\cite{roboref}, and Hu et al.\ inject false positives and false negatives
into the success detector of a hierarchical VLA at run time~\cite{hu2026orchestrating}. In this work, we inject
controlled errors into the training data of a policy that improves from its own rollouts.

\textbf{Self-improvement and reinforcement learning for VLAs.} Prior works have improved robot foundation models from
their own experience using learned success detectors or rewards~\cite{sefm,pistar06}, and have trained VLAs with
reinforcement learning against simulator success~\cite{simplevlarl}. These systems depend on the quality of their
success signal, but do not isolate how errors in that signal affect the result. Our experiments hold the policy,
rollouts, and training procedure fixed and vary only the errors in the success labels.

\textbf{Learning from noisy labels and imperfect verifiers.} Imitation from imperfect or noisy demonstrations has
been addressed by weighting demonstrations by confidence~\cite{wu2019imperfect,sasaki2021noisy}, and filtering
trajectories by return is a strong baseline in offline RL~\cite{chen2021dt}. Class-conditional label noise, in which
the two error rates differ, is well studied for classifiers~\cite{natarajan2013noisy}. In language models, policies
over-optimize imperfect reward models~\cite{gao2023overopt}, resampling is capped by verifier false
positives~\cite{flaws}, self-distillation gains track the filter's precision~\cite{zhou2026verifier}, and systematic
verifier errors can be more harmful than random ones~\cite{egashira2026delay}. We test whether these findings hold
when the verifier is a VLM, the learner is a robot policy, and the false positives are dominated by near-misses.

\section{Preliminaries}

We consider a policy $\pi_0$ that is rolled out on a set of tasks to produce trajectories $\tau$, each with a
ground-truth outcome $y(\tau) \in \{0, 1\}$. A grader assigns each trajectory a score $s(\tau)$, and trajectories with
$s(\tau) \geq t$ are kept. Self-improvement by filtered imitation then fine-tunes the policy on the kept set,
\begin{equation}
  \max_\theta \; \mathbb{E}_{\tau \sim \pi_0} \big[ \mathbb{1}[s(\tau) \geq t] \, \log \pi_\theta(\tau) \big],
\end{equation}
so that the policy imitates the distribution $p_t(\tau) \propto \pi_0(\tau)\,\mathbb{1}[s(\tau) \geq t]$. When the
grader is perfect, $p_t$ is the distribution of the policy's own successes. When it is not, let $q$ be the policy's
success rate, $\alpha$ the grader's false-positive rate (the fraction of failures kept), and $\beta$ its
false-negative rate (the fraction of successes rejected). The kept set is then a mixture of successes and failures,
in which failures make up a fraction
\begin{equation}
  \rho = \frac{\alpha (1 - q)}{\alpha (1 - q) + (1 - \beta)\, q} = 1 - \text{precision}.
  \label{eq:rho}
\end{equation}
This is the standard expression for one minus precision at prevalence $q$: $\rho$ is the grader's false discovery
rate on the policy's own rollouts. Unlike AUROC or accuracy, it depends on the policy's success rate as well as on the
grader. It makes the asymmetry between the two errors explicit. A false negative that is independent of the
trajectory only removes successes, which shrinks the kept set but leaves the distribution of kept successes
unchanged~\cite{elkan2008pu}.
A false positive changes the distribution itself, by adding failed trajectories to the imitation target. Equation
\eqref{eq:rho} also shows that raising the threshold reduces $\rho$ only if it reduces $\alpha$ faster than
$1 - \beta$. For a grader that cannot separate near-misses from successes, the two fall together, and a stricter
threshold discards data without making it cleaner.

\section{Setup}
\todo{Policy: \pizerofive{} LIBERO checkpoint. Rollouts: LIBERO-90, both cameras, state, actions, object and fixture
joint positions, goal predicates. Near-miss tagging: an object moved $>$3\,cm, a fixture joint moved, or part of the goal
holds. Graders: Robometer-4B (general and LIBERO fine-tuned), RoboReward-4B, Qwen yes/no prompts. Thresholds chosen on
a validation split by episode-name hash, metrics on test.}

\section{Experimental Evaluation}

Our experiments use the policy's own rollouts in LIBERO, where the simulator's success check serves as ground truth
for grading and for evaluating fine-tuned policies. They focus on the following questions:
\begin{enumerate}
  \item How often do current VLM graders pass failed attempts, and what kind of failures do they pass?
  \item How much in-domain data does a grader need to match or exceed a large general grader?
  \item Which grader error, false positives or false negatives, reduces the benefit of self-improvement?
  \item Is a real grader's effect on the policy explained by its error rates alone?
  \item Can a grader with high precision be obtained cheaply, and does it recover the benefit of self-improvement?
\end{enumerate}
Questions 1--4 diagnose why a general VLM grader fails to support self-improvement. Question 5, addressed in
Section~\ref{sec:fixes}, tests ways to make it work.

\subsection{How often do graders pass failures?}
\todo{AUROC with bootstrap CIs; false-positive rate at 90\% recall; near-miss vs.\ obvious failures; overlap of passed
failures between graders. Table I.}

\subsection{How much in-domain data does a grader need?}
\todo{Linear probes on frozen SigLIP and DINOv2 features vs.\ number of labels; seen vs.\ unseen tasks. Fig.~1.}

\subsection{Which grader error hurts self-improvement?}
\todo{Training pool of 19 partially solved tasks (570 rollouts). Datasets: Clean, FP~10/30/50, FN~10/30/50. LoRA
fine-tuning with identical settings; 285 held-out evaluation episodes per policy; paired bootstrap CIs; two seeds for
the main rows. Fig.~2: policy success vs.\ false-positive and false-negative rate.}

\subsection{Is a real grader explained by its error rates?}
\todo{Real grader vs.\ random errors matched per task to its rates; strict-threshold grader; relate results to
$\rho$ in Eq.~\eqref{eq:rho}. Contrast with~\cite{egashira2026delay}, where systematic false positives were worse than
random ones.}

\section{Making Self-Improvement Work}
\label{sec:fixes}

The results above identify what a grader needs in order to support self-improvement. Because rejecting successes is
nearly free, the grader does not need high recall. Because a real grader's effect is predicted by its error rates, the
share of failures in the data it keeps, $\rho$, can be estimated from a modest labelled set before any training. A grader
is therefore useful exactly when it can reach a low $\rho$ at some threshold while still keeping enough episodes on each
task. The general grader fails this test: on our training pool its kept set is \todo{30\%} failures at the balanced
threshold and \todo{24\%} at a strict one, because it cannot separate near-misses from successes. In this section we
test three ways to obtain a grader that passes it, each at a different cost.

\subsection{Candidate fixes}
\todo{(a) Domain adaptation: Robometer fine-tuned on LIBERO (AUROC 0.94 on the main pool vs.\ 0.84 for the general
model); an upper bound, since it was trained on LIBERO data. (b) A linear probe on frozen features trained on about 100
labelled rollouts (Section~V-B); the realistic option, since labelling a hundred of the robot's own attempts is cheap.
(c) Ranking within each task: keep each task's top-scoring episodes instead of applying one global threshold, since
grader scores are not calibrated across tasks (the strict global threshold lost \todo{5} of 19 tasks).}

\subsection{Predicting usefulness before training}
\todo{Table: for each grader and selection rule on the training pool, $\rho$, episodes kept, tasks covered. Computed
offline, no training. Identify which fixes reach $\rho \le$ \todo{10\%}.}

\subsection{Does a high-precision grader recover the gain?}
\todo{Fine-tune on the selections of the two or three most promising fixes (identical settings, 285 held-out episodes).
Fig.~3: policy success vs.\ $\rho$ for every run in the paper (Clean, FP sweep, real grader, matched, strict, fixes).
Claim to test: runs fall on one curve, and fixes with low $\rho$ land near Clean.}

\section{Discussion and Future Work}

\textbf{Toward recursive self-improvement.} Equation~\eqref{eq:rho} has a consequence beyond a single round. For a
fixed grader, $\rho$ falls as the policy's success rate $q$ rises, because failures make up a smaller part of what the
policy produces. A loop whose first round improves the policy therefore makes its own grading easier in the next round,
while a loop whose first round degrades the policy makes it harder. This suggests that recursive self-improvement with
a learned grader has a threshold: what matters is whether the first round starts above it, and the fixes in
Section~\ref{sec:fixes} are ways to start above it. \todo{State the threshold condition explicitly in terms of $\alpha$,
$\beta$, $q$; we test one round, so this is a prediction.}

\todo{Summary; implication for choosing and thresholding graders by precision; limitations (one policy, one benchmark,
simulation, one round of self-improvement); future work (multi-round loops that test the threshold prediction, graders
aimed at near-misses, real robots).}

\bibliographystyle{IEEEtran}
\bibliography{refs}

\end{document}
